\usepackage{amsmath, amssymb, amsthm}
\usepackage{mathtools, mathrsfs}

\usepackage{etoolbox, enumitem}

\setlist[enumerate,1]{
  label=\textnormal{(\roman*)},
  ref=\roman*,
  leftmargin=*,
  itemsep=2pt,
  topsep=3pt
}

\setlist[enumerate,2]{
  label=\textnormal{(\arabic*)},
  ref=\arabic*,
  leftmargin=*,
  itemsep=1pt,
  topsep=2pt
}

\setlist[itemize]{
  leftmargin=*,
  itemsep=2pt,
  topsep=3pt
}

\setlist[itemize]{
  % leftmargin=6pt,
  label=\(\vcenter{\hbox{\tiny\(\bullet\)}}\),
}

\makeatletter

\pretocmd{\@begintheorem}{%
  \begingroup
  \setlist[enumerate,1]{
    label=(\alph*),
    font=\normalfont
  }
  \setlist[enumerate,2]{
    label=(\roman*),
    font=\normalfont
  }
}{}{}

\apptocmd{\@endtheorem}{%
  \endgroup
}{}{}

\makeatother

% -------------------------------------------------
% Theorem Environments
% -------------------------------------------------

\newtheoremstyle{plain}
  {7pt}                         % Space above
  {4pt}                         % Space below
  {\itshape}                    % Body font
  {0pt}                         % Indentation
  {\bfseries}                   % Heading font
  {.}                           % Punctuation after heading
  {0.5em}                       % Space after heading
  {%
    \thmname{#1}%
    \thmnumber{ #2}%
    \thmnote{. (#3)}%
  }

\newtheoremstyle{upright}
  {7pt}                         % Space above
  {4pt}                         % Space below
  {\normalfont}                 % Body font
  {0pt}                         % Indentation
  {\bfseries}                   % Heading font
  {.}                           % Punctuation after heading
  {0.5em}                       % Space after heading
  {%
    \thmname{#1}%
    \thmnumber{ #2}%
    \thmnote{. #3}%
  }

\theoremstyle{plain}

\newtheorem{theorem}{Theorem}[section]
\newtheorem{definition}[theorem]{Definition}
\newtheorem{notation}[theorem]{Notation}
\newtheorem{lemma}[theorem]{Lemma}
\newtheorem{corollary}[theorem]{Corollary}
\newtheorem{proposition}[theorem]{Proposition}

\theoremstyle{upright}

\newtheorem{example}[theorem]{Example}
\newtheorem{remark}[theorem]{Remark}
\newtheorem{exercise}[theorem]{Exercise}

% -------------------------------------------------
% Bar Environments
% -------------------------------------------------
\usepackage{mdframed}

% proof: headful
\newmdenv[
  hidealllines=true,
  leftline=true,
  linewidth=1pt,
  leftmargin=0pt,
  rightmargin=0pt,
  innerleftmargin=8pt,
  innerrightmargin=0pt,
  innertopmargin=1pt,
  innerbottommargin=1pt,
  skipabove=6pt,
  skipbelow=7pt,
  splittopskip=\topskip,
  splitbottomskip=2pt
]{proofbar}

% roof: headless
\newmdenv[
  topline=false,
  bottomline=false,
  rightline=false,
  leftline=true,
  linewidth=1pt,
  linecolor=black,
  leftmargin=0pt,
  rightmargin=0pt,
  innerleftmargin=8pt,
  innerrightmargin=0pt,
  innertopmargin=2pt,
  innerbottommargin=2pt,
  skipabove=8pt,
  skipbelow=4pt,
  splittopskip=\topskip,
  splitbottomskip=2pt
]{roofbar}

\renewcommand{\qedsymbol}{\ensuremath{\blacktriangleleft}}

\renewenvironment{proof}[1][Proof]{%
  \par
  \pushQED{\qed}%
  \begin{proofbar}%
  \noindent\textit{#1.}\enspace\ignorespaces
}{%
  \popQED
  \par
  \end{proofbar}%
}

\newenvironment{solution}[1][Solution]
  {\begin{proof}[#1]}
  {\end{proof}}

\newenvironment{attempt}[1][Attempt]
  {\begin{proof}[#1]}
  {\end{proof}}

\newenvironment{answer}[1][Answer]
  {\begin{proof}[#1]}
  {\end{proof}}

\newenvironment{roof}
  {%
    \par
    \begin{roofbar}
    \ignorespaces
  }
  {%
    \par
    \end{roofbar}
  }

% -------------------------------------------------
% Headers
% -------------------------------------------------
\usepackage{fancyhdr}
\pagestyle{fancy}
\fancyhf{}

% Page number on the outer edge:
% even pages on the left, odd pages on the right.
\fancyhead[LE]{\thepage}
\fancyhead[RO]{\thepage}

% Even-page center: title.
\fancyhead[CE]{\nouppercase{\NotesTitle}}

% Odd-page center: current section.
\fancyhead[CO]{\nouppercase{\rightmark}}

\renewcommand{\headrulewidth}{0pt}
\renewcommand{\footrulewidth}{0pt}

% Use the chapter title until the first section begins.
\renewcommand{\chaptermark}[1]{%
  \markboth{#1}{#1}%
}

\renewcommand{\sectionmark}[1]{%
  \markright{#1}%
}

% Give chapter-opening pages the same header style.
\fancypagestyle{plain}{%
  \fancyhf{}
  \fancyhead[LE]{\thepage}
  \fancyhead[RO]{\thepage}
  \fancyhead[CE]{\nouppercase{\NotesTitle}}
  \fancyhead[CO]{\nouppercase{\rightmark}}
  \renewcommand{\headrulewidth}{0pt}
  \renewcommand{\footrulewidth}{0pt}
}


\newcommand{\NotesTitle}{Introduction to Algebra}
\newcommand{\NotesAuthor}{Tsz King Ho}

\newcommand{\makenotestitle}{%
  \begin{titlepage}
    \thispagestyle{empty}
    \centering
    \vspace*{\fill}

    {\Huge\bfseries \NotesTitle\par}

    \vspace{1.5cm}

    {\Large \NotesAuthor\par}

    \vspace*{\fill}
  \end{titlepage}
}

% -------------------------------------------------
% Macros
% -------------------------------------------------
\def\C{\mathbb{C}}
\def\E{\mathbb{E}} % expected value
\def\K{\mathbb{K}}
\def\R{\mathbb{R}}
\def\N{\mathbb{N}}
\def\P{\mathbb{P}} % probability
\def\Q{\mathbb{Q}}
\def\S{\mathbb{S}}
\def\T{\mathbb{T}}
\def\Z{\mathbb{Z}}

\def\O{\mathcal{O}}
\def\F{\mathcal{F}}
\def\L{\mathcal{L}}
\def\M{\mathcal{M}}

\def\cal{\mathcal}
\def\frak{\mathfrak}

\def\a{\alpha}
\def\b{\beta}
\def\c{\gamma}
\def\f{\varphi}
\def\eps{\varepsilon}
\def\lda{\lambda}
\def\Lda{\Lambda}
\def\sgm{\sigma}
\def\Sgm{\Sigma}

\newcommand{\pa}{\partial}
\newcommand{\cyc}{\mathrm{cyc}}
\newcommand{\ab}{\mathsf{ab}}
\newcommand{\PV}{\mathrm{P\!.V\!.}\,}

% basic
\DeclareMathOperator{\dom}{dom}
\DeclareMathOperator{\codom}{codom}
\DeclareMathOperator{\im}{im}
\DeclareMathOperator{\graph}{graph}
\DeclareMathOperator{\lcm}{lcm}

% calculus
\DeclareMathOperator{\degree}{deg}
\DeclareMathOperator{\proj}{proj}

\DeclareMathOperator{\dist}{dist}
\DeclareMathOperator{\diam}{diam}
\DeclareMathOperator{\supp}{supp}

\DeclareMathOperator*{\Res}{Res}
\DeclareMathOperator*{\argmax}{argmax}
\DeclareMathOperator*{\argmin}{argmin}
\DeclareMathOperator*{\esssup}{ess\,sup}
\DeclareMathOperator*{\essinf}{ess\,inf}

\DeclareMathOperator{\sgn}{sgn}
\DeclareMathOperator{\id}{id}
\DeclareMathOperator{\Id}{Id}
\renewcommand{\Re}{\operatorname{Re}}
\renewcommand{\Im}{\operatorname{Im}}
\DeclareMathOperator{\Arg}{Arg}
\DeclareMathOperator{\Sin}{Sin}
\DeclareMathOperator{\Cos}{Cos}
\DeclareMathOperator{\Tan}{Tan}
\DeclareMathOperator{\Log}{Log}
\DeclareMathOperator{\Exp}{Exp}

% topology
\DeclareMathOperator{\Int}{Int}
\DeclareMathOperator{\Cl}{Cl}
\DeclareMathOperator{\Bd}{Bd}
\DeclareMathOperator{\Frt}{Frt}

\newcommand{\std}{\mathsf{std}}

% probability
\DeclareMathOperator{\Var}{Var}
\DeclareMathOperator{\Cov}{Cov}
\DeclareMathOperator{\Corr}{Corr}
\DeclareMathOperator{\Bias}{Bias}

% analysis
\newcommand{\toweak}{\stackrel{w}{\rightharpoonup}}
\newcommand{\toast}{\stackrel{\ast}{\rightharpoonup}}

% linear algebra
\DeclareMathOperator{\rank}{rank}
\DeclareMathOperator{\nul}{nul}
\DeclareMathOperator{\rowsp}{rowsp}
\DeclareMathOperator{\colsp}{colsp}
\DeclareMathOperator{\spn}{span}
\DeclareMathOperator{\coker}{coker}
\DeclareMathOperator{\coim}{coim}
\DeclareMathOperator{\codim}{codim}
\DeclareMathOperator{\diag}{diag}
\DeclareMathOperator{\tr}{tr}

% algebra
\renewcommand{\lhd}{\triangleleft} % lhd uses mathbin and triangleleft uses mathrel, so the spacing is different
\renewcommand{\rhd}{\triangleright}
\DeclareMathOperator{\Map}{Map}
\DeclareMathOperator{\Hom}{Hom}
\DeclareMathOperator{\Obj}{Obj}
\DeclareMathOperator{\Aut}{Aut}
\DeclareMathOperator{\End}{End}
\DeclareMathOperator{\Ann}{Ann}
\DeclareMathOperator{\Inn}{Inn}
\DeclareMathOperator{\Gal}{Gal}
\DeclareMathOperator{\Spec}{Spec}
\DeclareMathOperator{\Frac}{Frac}
\DeclareMathOperator{\ch}{char}
\DeclareMathOperator*{\colim}{colim}

\DeclareMathOperator{\GL}{GL}
\DeclareMathOperator{\SL}{SL}
\DeclareMathOperator{\Sym}{Sym}

% restriction
\def\res{\mathord{\upharpoonright}}
\def\resmes{%
  \mathbin{
    \hspace{0.25ex}
    \vrule height 1.6ex width 0.08ex%
    \vrule height 0.08ex width 0.8ex%
    }%
}

% fourier
\newcommand\FTarrow{%
  \xrightarrow{\:\:\:\:\:\:\mathcal{F}\:\:\:\:\:\:}%
}
\newcommand\IFTarrow{%
  \xrightarrow{\:\:\:\mathcal{F}^{-1}\:\:\:}%
}

\newcommand{\wt}{\widetilde}
\newcommand{\wh}{\widehat}
\newcommand\widecheck[1]{%
    #1%
    \text{\llap{%
        \smash{\raisebox{\dimexpr 2\depth-0.5ex\relax}{%
            \scalebox{1}[-1]{\(\widehat{\phantom{#1}}\)}%
        }}%
    }\rule{0pt}{1.1em}}%
}
\newcommand\wc{\widecheck}

\def\<{\langle}
\def\>{\rangle}
\DeclarePairedDelimiter{\ceil}{\lceil}{\rceil}
\DeclarePairedDelimiter{\floor}{\lfloor}{\rfloor}

\usepackage{hyperref}